% Sources: W07-W21; L40-L47, L54-L60.
\boardsection{2}{估计}
\block{贝叶斯估计}
\[
\hat\theta_B:\quad\hat\theta_{\mathrm{Mo}},\quad
\hat\theta_{\mathrm{Me}},\quad\hat\theta_E.
\]
后验众数、中位数、期望值。
\[
\hat\theta_{\mathrm{Mo}}=\arg\max_\theta\pi(\theta\mid x),
\qquad
\int_0^{\hat\theta_{\mathrm{Me}}}\pi(\theta\mid x)\,d\theta=\frac12,
\qquad
\hat\theta_E=\E(\theta\mid x).
\]
后验分布不受损失函数影响：
\[
\pi(\theta\mid x,L)\propto p(x\mid\theta,L)\pi(\theta)
\quad\longrightarrow\quad\pi(\theta\mid x).
\]

\block{例2.2.1}
\[
X\mid\theta\sim N(\theta,100),\qquad\theta\sim N(100,225).
\]
\[
\text{后验均值}\quad
\frac{400+9x}{13}=\frac4{13}\times100+\frac9{13}x.
\]
\[
\text{后验方差}\quad\sigma_1^2
=\frac{100\times225}{100+225}=\frac{900}{13}.
\]
\[
\theta\mid X=x\sim N\left(\frac{400+9x}{13},\frac{900}{13}\right).
\]
当$x=115$，
\[
\frac{400+9\times115}{13}=\frac{1435}{13}\approx110.38,
\qquad\theta\mid X=115\sim N(110.38,69.23).
\]
\[
\text{后验标准差}=8.32.
\]

\block{例2.2.2}
\[
X\mid\theta\sim\Bin(n,\theta),\qquad
p(x\mid\theta)=\binom nx\theta^x(1-\theta)^{n-x}.
\]
选Beta先验：
\[
\theta\sim\Be(\alpha,\beta),\quad
\pi(\theta)=\frac1{B(\alpha,\beta)}
\theta^{\alpha-1}(1-\theta)^{\beta-1}.
\]
\[
B(\alpha,\beta)=\int_0^1t^{\alpha-1}(1-t)^{\beta-1}\,dt,\qquad
\E\theta=\frac\alpha{\alpha+\beta},\quad
\text{众数}=\frac{\alpha-1}{\alpha+\beta-2}.
\]
由贝叶斯公式，
\begin{align*}
\pi(\theta\mid x)&\propto p(x\mid\theta)\pi(\theta)\\
&\propto\theta^x(1-\theta)^{n-x}\theta^{\alpha-1}(1-\theta)^{\beta-1}\\
&\propto\theta^{\alpha+x-1}(1-\theta)^{\beta+n-x-1}.
\end{align*}
\[
\theta\mid x\sim\Be(\alpha+x,\beta+n-x),\quad
\pi(\theta\mid x)=
\frac{\theta^{\alpha+x-1}(1-\theta)^{\beta+n-x-1}}
{B(\alpha+x,\beta+n-x)}.
\]
后验众数。记$a=\alpha+x,\ b=\beta+n-x$。
\begin{align*}
\ell(\theta)&=(a-1)\log\theta+(b-1)\log(1-\theta),\\
\ell'(\theta)&=\frac{a-1}{\theta}-\frac{b-1}{1-\theta}=0,\\
(a-1)(1-\theta)&=\theta(b-1),\qquad
\theta=\frac{a-1}{a+b-2},\\
\ell''(\theta)&=-\frac{a-1}{\theta^2}-\frac{b-1}{(1-\theta)^2}<0.
\end{align*}
\[
\hat\theta_{\mathrm{Mo}}=\frac{\alpha+x-1}{\alpha+\beta+n-2}.
\]
下推后验均值估计：
\begin{align*}
\hat\theta_E
&=\E(\theta\mid x)
=\frac1{B(a,b)}\int_0^1\theta\,\theta^{a-1}(1-\theta)^{b-1}\,d\theta\\
&=\frac{B(a+1,b)}{B(a,b)}\\
&=\frac{\Gamma(a+1)\Gamma(b)}{\Gamma(a+b+1)}
\bigg/\frac{\Gamma(a)\Gamma(b)}{\Gamma(a+b)}
=\frac a{a+b}.
\end{align*}
利用$B(a,b)=\Gamma(a)\Gamma(b)/\Gamma(a+b)$，$\Gamma(a+1)=a\Gamma(a)$。
\[
\hat\theta_E=\frac{\alpha+x}{\alpha+\beta+n}.
\]
记$\alpha+\beta=s$，$\mu_0=\alpha/(\alpha+\beta)$，则
\[
\hat\theta_E=\frac{s}{s+n}\mu_0+\frac n{s+n}\frac xn.
\]
先验均值；样本频率。
\[
\theta\mid x=0\sim\Be(1,n+1).
\]
\begin{align*}
P(\theta>t\mid x=0)
&=\int_t^1(n+1)(1-\theta)^n\,d\theta
=(1-t)^{n+1}.
\end{align*}
\[
n=3,\ x=0,\ t=0.2:\quad0.8^{3+1}=0.8^4;
\qquad n=10,\ x=0:\quad0.8^{10+1}=0.8^{11}.
\]

\block{例2.4（2.2.3）}
\[
p(x\mid\theta)=e^{-(x-\theta)}\ind_{\{x\geq\theta\}},
\qquad X=\theta+Y,\quad Y\sim\Exp(1).
\]
柯西分布先验：
\[
\pi(\theta)=\frac1{\pi(1+\theta^2)},\qquad
\pi(\theta\mid x)=\frac{p(x\mid\theta)\pi(\theta)}{m(x)}.
\]
\begin{align*}
m(x)&=\int_{-\infty}^{+\infty}p(x\mid t)\pi(t)\,dt
=\frac{e^{-x}}\pi\int_{-\infty}^x\frac{e^t}{1+t^2}\,dt,\\
\pi(\theta\mid x)
&=\frac{e^{-x+\theta}}{\pi m(x)(1+\theta^2)}
\ind_{\{x\geq\theta\}}\\
&=\frac{e^\theta/(1+\theta^2)}
{\int_{-\infty}^xe^t/(1+t^2)\,dt}\ind_{\{x\geq\theta\}}.
\end{align*}
由于归一化常数不影响最大值的位置，故只需在$\theta\leq x$上最大化。
\[
g(\theta)=\frac{e^\theta}{1+\theta^2},\qquad
g'(\theta)=\frac{e^\theta(1+\theta^2)-2\theta e^\theta}{(1+\theta^2)^2}
=\frac{e^\theta(\theta-1)^2}{(1+\theta^2)^2}\geq0.
\]
\[
\hat\theta_{\mathrm{Mo}}=x=\hat\theta_{\mathrm{MLE}},\qquad
\E(\theta\mid x)=
\frac{\int_{-\infty}^x te^t/(1+t^2)\,dt}
{\int_{-\infty}^x e^t/(1+t^2)\,dt}<x.
\]

\block{估计误差}
\[
\MSE(\hat\theta\mid x)
=\Var(\theta\mid x)+\{\E(\theta\mid x)-\hat\theta\}^2.
\]
后验均方误差：
\[
\MSE(\hat\theta\mid x)
=\E\{(\theta-\hat\theta)^2\mid x\}
=\int(\theta-\hat\theta)^2\pi(\theta\mid x)\,d\theta.
\]
令$m=\E(\theta\mid x)$。给定$x$后，$m,\hat\theta$均为常数。
\begin{align*}
\theta-\hat\theta&=(\theta-m)+(m-\hat\theta),\\
(\theta-\hat\theta)^2
&=(\theta-m)^2+(m-\hat\theta)^2+2(\theta-m)(m-\hat\theta),\\
\MSE(\hat\theta\mid x)
&=\E\{(\theta-m)^2\mid x\}
+2(m-\hat\theta)\underbrace{\E(\theta-m\mid x)}_{=\,0}
+(m-\hat\theta)^2\\
&=\Var(\theta\mid x)+(m-\hat\theta)^2.
\end{align*}
特别地，
\[
\MSE(\hat\theta_E\mid x)=\Var(\theta\mid x),\qquad
\operatorname{SD}(\theta\mid x)=\sqrt{\Var(\theta\mid x)}.
\]
平方损失下，后验均值为最优贝叶斯估计！

\block{例2.5}
\[
P(X=x)=\theta(1-\theta)^{x-1},\qquad
P(X=3)=\theta(1-\theta)^2.
\]
$\theta$取$1/4,1/2,3/4$，
\[
P(\theta=1/4)=P(\theta=1/2)=P(\theta=3/4)=1/3.
\]
\begin{align*}
P(X=3\mid\theta=1/4)&=\tfrac14(\tfrac34)^2=\tfrac9{64},
&P(X=3,\theta=1/4)&=\tfrac9{192},\\
P(X=3\mid\theta=1/2)&=\tfrac12(\tfrac12)^2=\tfrac8{64},
&P(X=3,\theta=1/2)&=\tfrac8{192},\\
P(X=3\mid\theta=3/4)&=\tfrac34(\tfrac14)^2=\tfrac3{64},
&P(X=3,\theta=3/4)&=\tfrac3{192}.
\end{align*}
由全概率公式，
\[
P(X=3)=\frac9{192}+\frac8{192}+\frac3{192}
=\frac{20}{192}=\frac5{48}.
\]
由贝叶斯公式，
\[
P(\theta=t\mid X=3)=\frac{P(X=3,\theta=t)}{P(X=3)}.
\]
\begin{center}
\begin{tabular}{c|ccc}
$\theta$&$1/4$&$1/2$&$3/4$\\\hline
$P(\theta\mid X=3)$&$9/20$&$8/20$&$3/20$
\end{tabular}
\end{center}
\[
\hat\theta_{\mathrm{Mo}}=\frac14,\qquad
\hat\theta_E=\E(\theta\mid X=3)
=\frac14\frac9{20}+\frac12\frac8{20}+\frac34\frac3{20}
=\frac{17}{40}=0.425.
\]
\[
\E(\theta^2\mid X=3)
=\left(\frac14\right)^2\frac9{20}
+\left(\frac12\right)^2\frac8{20}
+\left(\frac34\right)^2\frac3{20}=\frac{17}{80}.
\]
\[
\Var(\theta\mid X=3)=\frac{17}{80}-\left(\frac{17}{40}\right)^2
=\frac{51}{1600}.
\]
\[
\MSE(\hat\theta_{\mathrm{Mo}}\mid X=3)
=\frac{51}{1600}+\left(\frac{17}{40}-\frac14\right)^2
=\frac1{16},\qquad
\sqrt{\MSE(\hat\theta_{\mathrm{Mo}}\mid X=3)}=\frac14.
\]
\[
\E\{(\theta-\tfrac14)^2\mid X=3\}
=0^2\frac9{20}+\left(\frac14\right)^2\frac8{20}
+\left(\frac12\right)^2\frac3{20}=\frac1{16}.
\]

\block{例2.2.5}
\[
X\mid\theta\sim\Bin(n,\theta),\qquad\theta\sim\Be(\alpha,\beta),
\]
\[
p(x\mid\theta)\propto\theta^x(1-\theta)^{n-x},\qquad
\pi(\theta)\propto\theta^{\alpha-1}(1-\theta)^{\beta-1}.
\]
\[
\pi(\theta\mid x)\propto\theta^{\alpha+x-1}(1-\theta)^{\beta+n-x-1},
\qquad\theta\mid x\sim\Be(\alpha+x,\beta+n-x).
\]
\[
\E(\theta\mid x)=\frac{\alpha+x}{\alpha+\beta+n},\qquad
\Var(\theta\mid x)=
\frac{(\alpha+x)(\beta+n-x)}
{(\alpha+\beta+n)^2(\alpha+\beta+n+1)}.
\]
特别地，若$\alpha=\beta=1$，即$\theta\sim U(0,1)$，有
\[
\theta\mid x\sim\Be(x+1,n-x+1),\qquad
\hat\theta_E=\frac{x+1}{n+2}
=\frac n{n+2}\frac xn+\frac2{n+2}\frac12.
\]
\[
\frac d{d\theta}\log\pi(\theta\mid x)
=\frac x\theta-\frac{n-x}{1-\theta}=0,\qquad
\hat\theta_{\mathrm{Mo}}=\frac xn.
\]
\[
\MSE(\hat\theta_E\mid x)=\Var(\theta\mid x)
=\frac{(x+1)(n-x+1)}{(n+2)^2(n+3)}.
\]
% The missing square in this next denominator is visible in L46-L47;
% retained literally, with the discrepancy recorded in transcription_checks.md.
\[
\MSE(\hat\theta_{\mathrm{Mo}}\mid x)
=\frac{(x+1)(n-x+1)}{(n+2)(n+3)}
+\left(\frac{x+1}{n+2}-\frac xn\right)^2.
\]
\[
\MSE(\hat\theta_{\mathrm{Mo}}\mid x)-\MSE(\hat\theta_E\mid x)
=\frac{(n-2x)^2}{n^2(n+2)^2}\geq0.
\]
如$n=3,x=0$，$\theta\mid x\sim\Be(1,4)$。
\[
\hat\theta_E=\frac15,\quad\hat\theta_{\mathrm{Mo}}=0,\qquad
\MSE(\hat\theta_E\mid x)=\frac{1\times4}{5^2\times6}=\frac2{75},
\]
\[
\MSE(\hat\theta_{\mathrm{Mo}}\mid x)
=\frac2{75}+\left(\frac15\right)^2=\frac1{15}.
\]
